\subsection{A general first-unit bound and its ramification boundary}

The same principal-part argument has a useful general form.
Its prime-to-$p$ scope must be retained: outside the special
weight, the first unit index could be $p$ itself.

\begin{proposition}[KT-COHEN-GENERAL-FIRST-UNIT-BOUND]
Let $p\geq29$ be prime, and let $b$ be odd with
\[
 3<b\leq p-2,\qquad b\ne(p-1)/2,
 \qquad C=-B_{2b}/(2b)\in\mathbf Z_p^\times.
 \tag{CG.1}
\]
Then there exists an odd negative fundamental discriminant
$D=-d$ such that
\[
 3\leq d\leq2b-7<2p,
 \qquad p\nmid h(D),\qquad
 \beta_b(\chi_D)\in\mathbf Z_p^\times.
 \tag{CG.2}
\]
If either $p\equiv1\pmod4$ or $2b-7<p$, the discriminant
in (CG.2) is automatically prime to $p$.  In these cases
the two generalized $B_1$ units and the all-weight transport
of KT-JACOBI-TWO-UNIT-FIELD apply to this same $D$.
For $b=3$, the prime-to-$p$ choice $D=-3$ has
$h(D)=1$ and $\beta_3(\chi_D)=2/9$.
\end{proposition}

\begin{proof}
Use the notation and normalizations of
KT-COHEN-WEIGHT-TWO-BOUND, but set
\[
                  m=2b-3,\qquad F_b=H_b/\theta^m.
 \tag{CG.3}
\]
The difference of theta-multiplier exponents is again
$2b+1-m=4$, so this is weight two with trivial character.
It is meromorphic only at the cusp $1/2$.
The complete coefficient-integrality proof in
KT-JACOBI-TWO-UNIT-FIELD applies under (CG.1), including
at indices divisible by $p$.

Since $H_b=C+O(q^3)$ and $\theta=1+2q+O(q^4)$,
the ordinary binomial expansion gives
\[
 F_b=C-2mCq+2m(m+1)Cq^2+O(q^3).
 \tag{CG.4}
\]
The cusp calculation (CW.10)--(CW.11) has not changed.
Its numerator terms are $p$-adic units times
$H(b,n)q^{n/4}$ for $n\equiv3\pmod4$;
the denominator in (CG.3) is $2^m q^{m/4}$ times
an integral power-series unit.  Here $m\equiv3\pmod4$.
Consequently, if every $H(b,n)$ with $0<n<m$ and
$n\equiv3\pmod4$ vanished modulo $p$, the principal
part of $F_b$ at $1/2$ would be divisible by $p$.

The explicit forms $W_j=G^{j+1}/\Theta^j$ in
(CW.12)--(CW.13) remove that principal part, using only
multiples divisible by $p$.  This would produce a
holomorphic characteristic-zero weight-two form, integral
at $p$, with the same three displayed coefficients modulo
$p$ as (CG.4).  The same coefficient-field Sturm bound
one applies.
Both basis forms $\Theta$ and $G$ have their $q^2$
coefficient equal to $24$ times their constant coefficient.
Their constant/$q$ coefficient determinant is $16$, so
the Sturm theorem implies this relation for every
holomorphic weight-two form modulo $p$ at level $4$.
For (CG.4) the required relation would be
\[
 \begin{aligned}
 0&=\bigl(2m(m+1)-24\bigr)C\\
  &=2(m+4)(m-3)C
    =4(2b+1)(b-3)C\pmod p.
 \end{aligned}
 \tag{CG.5}
\]
This is impossible.  The factors $4$ and $C$ are units;
$0<b-3<p$; and $0<2b+1<2p$, whose only possible
multiple of $p$ would require $b=(p-1)/2$, explicitly
excluded in (CG.1).

There is therefore a unit odd-index coefficient at some
$n<m$.  The minimal such index is fundamental, by exactly
the integral divisor-factor argument (CW.16).
Since both $m$ and $n$ are $3$ modulo $4$, it satisfies
$d=n\leq m-4=2b-7$.
The carried elementary class-number bound applies since
$d<2p<4p$, so $h(D)<p$.

The only positive multiple of $p$ that could lie below
$2p$ is $p$.  It cannot equal $d$ if $p\equiv1\pmod4$,
because $d\equiv3\pmod4$; it cannot occur if $d<p$
either.  This proves the asserted prime-to-$p$ cases.
There the existing character transport has its required
conductor hypothesis.  The separate $b=3$ assertion is
the stated exact elementary value and class number.
\end{proof}

\paragraph{What the general bound does not imply.}
For $p\equiv3\pmod4$ and $2b-7\geq p$, this proof
does not rule out $d=p$.  Integrality of a generalized
Bernoulli value at conductor $p$ does not imply its
divisibility by $p$.  No claim that (CG.2) always satisfies
the unramified auxiliary-field condition is made in that
range.  The earlier heat construction supplies its own
prime-to-$p$ field with its earlier bound; combining the
two arguments does not assign the sharper bound to that
field without a further step.

At the special weight $b=(p+3)/2$ with $p\equiv3\pmod4$,
one has $2b-7=p-4$, recovering (CW.1).
Neither the general nor the special bound proves a
same-discriminant unit for the additional twisted
Bernoulli value.  This is an auxiliary first-unit bound,
not a new vanishing statement for the original rational
class-group component.
